% GENERATED FILE. Edit canonical lessons, then run npm run generate:books.
% canonical-source-hash: fnv1a64:1497cbe32f660f2e
% canonical-lessons: RU-C222-yeda, RU-C222-blyudo, RU-C222-menyu, RU-C222-bulka, RU-C222-kefir

\chapter{Food, Dish, Menu, Bun, Kefir}
\label{ch:ru-food-dish-menu-bun-kefir}
\hlchaptermodality{222}

\begin{chapteropening}
\textbf{By the end of this chapter:} \emph{I can say food, a dish, a menu, a bun and kefir.}

Complete the last lesson of chapter 222: I can say food, a dish, a menu, a bun and kefir.
\end{chapteropening}

\section[\texorpdfstring{\ru{еда} (yedá)}{еда (yedá)}]{\ru{еда} (yedá) — food}

\label{lesson:RU-C222-yeda}



\begin{quote}
Before the new one: say the Russian for a vase, then the Russian for a painting.
\end{quote}

\subsection*{You'll want to know: \ru{еда}}
\textbf{\ru{еда}} (\emph{yedá}) — \textquotedblleft{}food\textquotedblright{}.

\begin{grammarlens}[title={What you've built}]
One more word for this chapter.
\end{grammarlens}

\subsection*{Your turn}
\begin{itemize}
\raggedright
  \item \emph{Say it:} \emph{yedá}
  \item \emph{Say it:} \emph{yedá}, once more
  \item \emph{Say it:} say \emph{yedá}
  \item \emph{Recall:} say the Russian for a vase, then the Russian for a painting, then say \emph{yedá} again
\end{itemize}

\begin{tcolorbox}[breakable,colback=teal!4,colframe=teal!35!black,title={Before you move on}]
What does \ru{еда} mean? (\textquotedblleft{}Food\textquotedblright{}.) Say it once more.
\end{tcolorbox}
\section[\texorpdfstring{\ru{блюдо} (blyúdo)}{блюдо (blyúdo)}]{\ru{блюдо} (blyúdo) — a dish}

\label{lesson:RU-C222-blyudo}



\begin{quote}
Before the new one: say the Russian for a painting, then the Russian for food.
\end{quote}

\subsection*{You'll want to know: \ru{блюдо}}
\textbf{\ru{блюдо}} (\emph{blyúdo}) — \textquotedblleft{}a dish\textquotedblright{}.

\begin{grammarlens}[title={What you've built}]
One more word for this chapter.
\end{grammarlens}

\subsection*{Your turn}
\begin{itemize}
\raggedright
  \item \emph{Say it:} \emph{blyúdo}
  \item \emph{Say it:} \emph{blyúdo}, once more
  \item \emph{Say it:} say \emph{blyúdo}
  \item \emph{Recall:} say the Russian for a painting, then the Russian for food, then say \emph{blyúdo} again
\end{itemize}

\begin{tcolorbox}[breakable,colback=teal!4,colframe=teal!35!black,title={Before you move on}]
What does \ru{блюдо} mean? (\textquotedblleft{}A dish\textquotedblright{}.) Say it once more.
\end{tcolorbox}
\section[\texorpdfstring{\ru{меню} (menyú)}{меню (menyú)}]{\ru{меню} (menyú) — a menu}

\label{lesson:RU-C222-menyu}



\begin{quote}
Before the new one: say the Russian for food, then the Russian for a dish.
\end{quote}

\subsection*{You'll want to know: \ru{меню}}
\textbf{\ru{меню}} (\emph{menyú}) — \textquotedblleft{}a menu\textquotedblright{}.

\begin{grammarlens}[title={What you've built}]
One more word for this chapter.
\end{grammarlens}

\subsection*{Your turn}
\begin{itemize}
\raggedright
  \item \emph{Say it:} \emph{menyú}
  \item \emph{Say it:} \emph{menyú}, once more
  \item \emph{Say it:} say \emph{menyú}
  \item \emph{Recall:} say the Russian for food, then the Russian for a dish, then say \emph{menyú} again
\end{itemize}

\begin{tcolorbox}[breakable,colback=teal!4,colframe=teal!35!black,title={Before you move on}]
What does \ru{меню} mean? (\textquotedblleft{}A menu\textquotedblright{}.) Say it once more.
\end{tcolorbox}
\section[\texorpdfstring{\ru{булка} (búlka)}{булка (búlka)}]{\ru{булка} (búlka) — a bun, a roll}

\label{lesson:RU-C222-bulka}



\begin{quote}
Before the new one: say the Russian for a dish, then the Russian for a menu.
\end{quote}

\subsection*{You'll want to know: \ru{булка}}
\textbf{\ru{булка}} (\emph{búlka}) — \textquotedblleft{}a bun, a roll\textquotedblright{}.

\begin{grammarlens}[title={What you've built}]
One more word for this chapter.
\end{grammarlens}

\subsection*{Your turn}
\begin{itemize}
\raggedright
  \item \emph{Say it:} \emph{búlka}
  \item \emph{Say it:} \emph{búlka}, once more
  \item \emph{Say it:} say \emph{búlka}
  \item \emph{Recall:} say the Russian for a dish, then the Russian for a menu, then say \emph{búlka} again
\end{itemize}

\begin{tcolorbox}[breakable,colback=teal!4,colframe=teal!35!black,title={Before you move on}]
What does \ru{булка} mean? (\textquotedblleft{}A bun, a roll\textquotedblright{}.) Say it once more.
\end{tcolorbox}
\section[\texorpdfstring{\ru{кефир} (kefír)}{кефир (kefír)}]{\ru{кефир} (kefír) — kefir}

\label{lesson:RU-C222-kefir}



\begin{quote}
Before the new one: say the Russian for a menu, then the Russian for a bun.
\end{quote}

\subsection*{You'll want to know: \ru{кефир}}
\textbf{\ru{кефир}} (\emph{kefír}) — \textquotedblleft{}kefir\textquotedblright{}.

\begin{grammarlens}[title={What you've built}]
One more word for this chapter.
\end{grammarlens}

\subsection*{Your turn}
\begin{itemize}
\raggedright
  \item \emph{Say it:} \emph{kefír}
  \item \emph{Say it:} \emph{kefír}, once more
  \item \emph{Say it:} say \emph{kefír}
  \item \emph{Recall:} say the Russian for a menu, then the Russian for a bun, then say \emph{kefír} again
\end{itemize}

\begin{tcolorbox}[breakable,colback=teal!4,colframe=teal!35!black,title={Before you move on}]
What does \ru{кефир} mean? (\textquotedblleft{}Kefir\textquotedblright{}.) Say it once more.
\end{tcolorbox}
